\section{2018级工科数学分析下 B卷}

\subsection{资料来源}
\begin{itemize}
  \item 试题：2018 级工科数学分析（二）B 卷，已导出页面图校正公式。
  \item 答案：B 卷附解答文件，已导出页面图校正公式。
\end{itemize}

\subsection{试题}

\paragraph*{试卷信息}
诚信应考，考试作弊将带来严重后果！

华南理工大学本科生期末考试《工科数学分析（二）》B卷，2018--2019学年第二学期。

注意事项：开考前请将密封线内各项信息填写清楚；所有答案请直接答在试卷上；考试形式：闭卷；本试卷共5大题，满分100分，考试时间120分钟。评阅教师请在试卷袋上评阅栏签名。

\paragraph*{一、填空题（每小题3分，共15分）}
\begin{enumerate}
  \item 设函数 \(u(x,y,z)=\sqrt{x^2+y^2+z^2}\)，则 \(\operatorname{div}(\operatorname{grad}u)=\) \underline{\hspace{5cm}}。
  \item 交换积分次序，则
  \[
    \int_0^2\mathrm{d}y\int_{y^2}^{4}f(x,y)\,\mathrm{d}x=
  \]
  \underline{\hspace{5cm}}。
  \item 设曲线 \(L:\dfrac{x^2}{a^2}+\dfrac{y^2}{b^2}=1\)，取顺时针方向，则
  \[
    \oint_L(x-y)\,\mathrm{d}x+(x+y)\,\mathrm{d}y=
  \]
  \underline{\hspace{5cm}}。
  \item 若级数 \(\displaystyle \sum_{n=1}^{\infty}\frac{(-1)^n(n+2)}{n^p}\) 条件收敛，则常数 \(p\) 的取值范围为 \underline{\hspace{5cm}}。
  \item 设
  \[
    f(x)=
    \begin{cases}
      -1, & -1<x\le0,\\
      x^2+2, & 0<x\le1,
    \end{cases}
  \]
  设 \(S(x)\) 为 \(f(x)\) 展成以 \(2\) 为周期的傅里叶级数的和函数。
  则 \(S(0)=\) \underline{\hspace{5cm}}。
\end{enumerate}

\paragraph*{二、计算题（每小题9分，共36分）}
\begin{enumerate}
  \item 设 \(u=f(x+y+z,x^2+y^2+z^2)\)，其中函数 \(f\) 有二阶连续的偏导数，求
  \[
    \frac{\partial^2u}{\partial x^2}+
    \frac{\partial^2u}{\partial y^2}+
    \frac{\partial^2u}{\partial z^2}.
  \]

  \item 计算三重积分
  \[
    \iiint_\Omega\left(\frac x2-y-z\right)^2\,\mathrm{d}x\mathrm{d}y\mathrm{d}z,
  \]
  其中 \(\Omega\) 是曲面 \(\dfrac{x^2}{4}+y^2+z^2=1\) 所围成的区域。

  \item 求抛物面 \(\Sigma:z=2-(x^2+y^2)\) 的质量（\(z\ge0\) 的部分），其密度函数为
  \(\rho(x,y,z)=x^2+y^2\)。

  \item 求微分方程
  \[
    x(2+y)\,\mathrm{d}x+y(1-x)\,\mathrm{d}y=0
  \]
  的通解。
\end{enumerate}

\paragraph*{三、解答题（每小题10分，共30分）}
\begin{enumerate}
  \item 设曲线积分
  \[
    \int_L\left(\cos x-\varphi(x)\right)\frac yx\,\mathrm{d}x+\varphi(x)\,\mathrm{d}y
  \]
  与积分路径无关，其中 \(\varphi(x)\) 有一阶连续导数且
  \(\varphi\left(\dfrac{\pi}{2}\right)=0\)，求 \(\varphi(x)\)，并计算曲线积分
  \[
    \int_{(1,0)}^{(\pi,\pi)}
    \left(\cos x-\varphi(x)\right)\frac yx\,\mathrm{d}x+\varphi(x)\,\mathrm{d}y.
  \]

  \item 计算
  \[
    I=\oiint_\Sigma
    \frac{x\,\mathrm{d}y\mathrm{d}z+y\,\mathrm{d}z\mathrm{d}x+z\,\mathrm{d}x\mathrm{d}y}
    {(x^2+y^2+z^2)^{3/2}},
  \]
  其中 \(\Sigma\) 为有界单连通区域的闭边界曲面，取外侧。

  \item 求幂级数 \(\displaystyle \sum_{n=0}^{\infty}\frac{3n+1}{n!}x^{3n}\) 的收敛域及和函数。
\end{enumerate}

\paragraph*{四、证明题（本题10分）}
证明函数项级数
\[
  \sum_{n=1}^{\infty}\frac{\cos(nx)}{n^2}
\]
在 \((-\infty,+\infty)\) 上一致收敛。

\paragraph*{五、应用题（本题9分）}
欲造一无盖的长方体容器，已知底部造价为每平方米3元，侧面造价均为每平方米1元，现想用36元造一个容积最大的容器，求它的尺寸。

\subsection{答案与解析}

\paragraph*{一、填空题}
\begin{enumerate}
  \item 当 \((x,y,z)\ne(0,0,0)\) 时，
  \[
    \operatorname{div}(\operatorname{grad}u)=\Delta\sqrt{x^2+y^2+z^2}
    =\frac{2}{\sqrt{x^2+y^2+z^2}}.
  \]
  \item 区域为 \(0\le y\le2,\ y^2\le x\le4\)，换序得
  \[
    \int_0^2\mathrm{d}y\int_{y^2}^{4}f(x,y)\,\mathrm{d}x
    =\int_0^4\mathrm{d}x\int_0^{\sqrt x}f(x,y)\,\mathrm{d}y.
  \]
  \item 由 Green 公式，逆时针时
  \[
    \oint_L(x-y)\,\mathrm{d}x+(x+y)\,\mathrm{d}y
    =\iint_D 2\,\mathrm{d}A=2\pi ab.
  \]
  题中取顺时针方向，故结果为
  \[
    -2\pi ab.
  \]
  \item 交错收敛要求 \((n+2)/n^p\to0\)，即 \(p>1\)；绝对收敛要求 \(\sum n^{1-p}\) 收敛，即 \(p>2\)。故条件收敛范围为
  \[
    1<p\le2.
  \]
  \item 傅里叶级数在跳跃点取左右极限平均值：
  \[
    S(0)=\frac{f(0-0)+f(0+0)}2=\frac{-1+2}{2}=\frac12.
  \]
\end{enumerate}

\paragraph*{二、计算题}
\begin{enumerate}
  \item 记
  \[
    s=x+y+z,\qquad r=x^2+y^2+z^2.
  \]
  有
  \[
    u_x=f_1+2xf_2,
  \]
  从而
  \[
    u_{xx}=f_{11}+4xf_{12}+4x^2f_{22}+2f_2.
  \]
  同理可得 \(u_{yy}\)、\(u_{zz}\)，相加得到
  \[
    u_{xx}+u_{yy}+u_{zz}
    =3f_{11}+4(x+y+z)f_{12}+4(x^2+y^2+z^2)f_{22}+6f_2.
  \]
  各偏导数均在 \((x+y+z,x^2+y^2+z^2)\) 处取值。

  \item 展开平方后由椭球区域对称性，交叉项积分为 \(0\)，所以
  \[
    \iiint_\Omega\left(\frac x2-y-z\right)^2\,\mathrm{d}V
    =\iiint_\Omega\left(\frac{x^2}{4}+y^2+z^2\right)\,\mathrm{d}V.
  \]
  作广义球坐标变换
  \[
    x=2r\sin\varphi\cos\theta,\quad
    y=r\sin\varphi\sin\theta,\quad
    z=r\cos\varphi,\quad J=2r^2\sin\varphi.
  \]
  因此
  \[
    I=2\int_0^{2\pi}\int_0^\pi\int_0^1r^4\sin\varphi\,\mathrm{d}r\,\mathrm{d}\varphi\,\mathrm{d}\theta
    =\frac{8\pi}{5}.
  \]

  \item 投影区域为 \(D:x^2+y^2\le2\)，且
  \[
    \mathrm{d}S=\sqrt{1+4x^2+4y^2}\,\mathrm{d}x\mathrm{d}y.
  \]
  因而质量
  \[
    M=\iint_D(x^2+y^2)\sqrt{1+4x^2+4y^2}\,\mathrm{d}x\mathrm{d}y
    =2\pi\int_0^{\sqrt2}r^3\sqrt{1+4r^2}\,\mathrm{d}r
    =\frac{149\pi}{30}.
  \]

  \item 原方程分离为
  \[
    \frac{y\,\mathrm{d}y}{2+y}=\frac{x\,\mathrm{d}x}{x-1}.
  \]
  两边积分得
  \[
    y-2\ln|y+2|=x+\ln|x-1|+C,
  \]
  等价地可写为
  \[
    (x-1)(y+2)^2=C\mathrm{e}^{y-x}.
  \]
\end{enumerate}

\paragraph*{三、解答题}
\begin{enumerate}
  \item 取右半平面作为单连通区域。令
  \[
    P=\left(\cos x-\varphi(x)\right)\frac yx,\qquad Q=\varphi(x).
  \]
  路径无关要求 \(P_y=Q_x\)，即
  \[
    \varphi'(x)+\frac1x\varphi(x)=\frac{\cos x}{x}.
  \]
  因此
  \[
    (x\varphi(x))'=\cos x,\qquad
    x\varphi(x)=\sin x+C.
  \]
  由 \(\varphi(\pi/2)=0\)，得 \(C=-1\)，故
  \[
    \varphi(x)=\frac{\sin x-1}{x}.
  \]
  势函数可取 \(F(x,y)=y\varphi(x)\)，于是
  \[
    \int_{(1,0)}^{(\pi,\pi)}P\,\mathrm{d}x+Q\,\mathrm{d}y
    =F(\pi,\pi)-F(1,0)=-1.
  \]

  \item 记 \(\Sigma\) 围成的区域为 \(\Omega\)。若原点不在 \(\Omega\) 内，则向量场
  \[
    \mathbf F=\frac{(x,y,z)}{(x^2+y^2+z^2)^{3/2}}
  \]
  在 \(\Omega\) 内无奇点，且散度为 \(0\)，由 Gauss 公式得
  \[
    I=0.
  \]
  若原点在 \(\Omega\) 内，则挖去小球 \(x^2+y^2+z^2=\varepsilon^2\)。外边界和内侧小球面合用 Gauss 公式，得
  \[
    I=\frac1{\varepsilon^3}\oiint_{x^2+y^2+z^2=\varepsilon^2}
    x\,\mathrm{d}y\mathrm{d}z+y\,\mathrm{d}z\mathrm{d}x+z\,\mathrm{d}x\mathrm{d}y
    =4\pi.
  \]

  \item 记
  \[
    s(x)=\sum_{n=0}^{\infty}\frac{3n+1}{n!}x^{3n}.
  \]
  因为
  \[
    \lim_{n\to\infty}\frac{a_{n+1}}{a_n}=0,
  \]
  收敛半径为 \(+\infty\)，收敛域为 \((-\infty,+\infty)\)。又
  \[
    s(x)=3\sum_{n=0}^{\infty}\frac{n(x^3)^n}{n!}
    +\sum_{n=0}^{\infty}\frac{(x^3)^n}{n!}
    =(3x^3+1)\mathrm{e}^{x^3}.
  \]
\end{enumerate}

\paragraph*{四、证明题}
对任意 \(x\in(-\infty,+\infty)\)，有
\[
  \left|\frac{\cos(nx)}{n^2}\right|\le\frac1{n^2}.
\]
又 \(\sum_{n=1}^{\infty}1/n^2\) 收敛，由 Weierstrass 判别法，原函数项级数在 \((-\infty,+\infty)\) 上一致收敛。

\paragraph*{五、应用题}
设长方体的长、宽、高分别为 \(x,y,z\) 米。费用约束为
\[
  3xy+2xz+2yz=36.
\]
最大化 \(V=xyz\)。令
\[
  F=\ln x+\ln y+\ln z+\lambda(3xy+2xz+2yz-36).
\]
由
\[
  \frac1x+3\lambda y+2\lambda z=0,\quad
  \frac1y+3\lambda x+2\lambda z=0,\quad
  \frac1z+2\lambda x+2\lambda y=0
\]
解得唯一临界点
\[
  x=2,\qquad y=2,\qquad z=3.
\]
实际问题中最大体积存在，故长、宽、高分别为 \(2\) 米、\(2\) 米、\(3\) 米。
